\subsection{From the cubical generator to the actual polarising frame}
\label{subsec:ftb-cubical-complex-frame}

Keep the actual exponential \(e:V=\mathbb C^g\to P(\mathbb C)\) and
lattice \(\Lambda\) of
Proposition~\ref{prop:ftb-actual-complex-exponential}.
Let \(L\) be the polarising line bundle, with its prescribed
rigidification at the origin.  The algebraic input in this step is the
source's \(\mathrm{CubeSt}\), applied to this rigidified \(L\).
It supplies a nowhere vanishing algebraic generator \(t\) of
\[
 \Theta(L)=\bigotimes_{I\subset\{0,1,2\}}
              s_I^*L^{\,\sigma_I},\qquad
 s_I(x)=\sum_{a\in I}x_a,\qquad \sigma_I=(-1)^{3-|I|}.        \tag{CF1}
\]
Only this particular bundle and generator are used.

Write \(U_i\) for the algebraic charts of \(L\), and \(g_{ij}\) for their
transition units, in the convention that section coordinates satisfy
\(s_i=g_{ij}s_j\).  Define
\[
 W_i=e^{-1}U_i,\qquad c_{ij}(z)=g_{ij}(e(z)).
\]
For a choice \(j=(j_I)_{I\subset\{0,1,2\}}\), let \(T_j(x)\) be the value
of the \(j\)-coordinate of \(t\) at
\((e(x_0),e(x_1),e(x_2))\).  The charts \(W_i\) cover \(V\).
The \(c_{ij}\) are nonzero holomorphic transition functions and
\[
 c_{ij}c_{jk}=c_{ik},\qquad
 T_j(x)=\prod_I c_{j_Ij'_I}(s_Ix)^{\sigma_I}\,T_{j'}(x).    \tag{CF2}
\]
Each \(T_j\) is holomorphic and nonzero on the open set where all the
\(s_Ix\) lie in \(W_{j_I}\).  Moreover \(W_i,c_{ij}\) are
\(\Lambda\)-periodic, and \(T_j\) is periodic under independent
\(\Lambda\)-translations of the three variables.
These assertions follow by evaluating the algebraic transition and
generator identities: the partial sum of the three points is
\(e(s_Ix)\), and the cube transition unit evaluates to the product in
(CF2).  Thus this is the pulled-back algebraic cube, not a separately
chosen analytic cubical object.

\begin{srlemma}[The source's local third-difference construction]
\label{lem:ftb-cube-local-third-difference}
The preceding actual data admit a radius \(r>0\), nonvanishing
holomorphic functions \(b_i\) on \(W_i\cap B(0,r)\), and a scalar
\(\kappa\ne0\), such that
\[
 b_i=c_{ij}b_j,\qquad
 T_j(x)=\kappa\prod_I b_{j_I}(s_Ix)^{\sigma_I}                \tag{CF3}
\]
for all sufficiently small \(x\) in the relevant cube chart.
\end{srlemma}

\begin{proof}
Choose \(i_0\) with \(0\in W_{i_0}\), and put
\(F(x)=T_{(i_0,\ldots,i_0)}(x)\).
We first give the actual face and doubling argument.
On a face \(x_a=0\), every subset term is paired with the subset
obtained by inserting or deleting \(a\).  The transition products in
(CF2) cancel in pairs.  Hence the corresponding local \(T_j\)'s glue to
a single entire function on that face, periodic under a real-spanning
lattice.  It is constant: choose a real lattice basis, reduce to a
compact fundamental parallelepiped, and apply Liouville's theorem.
Its value at zero is \(F(0)\).  Thus \(F(x)=F(0)\) on each small face.

There is one further glued entire periodic quotient.  Its local
expression is
\[
 \frac{T_{j^+}(2x)
       \prod_I T_{\operatorname{diag}(j_I,j^-_I,j^+_I,i_0)}
                    (s_Ix,s_Ix,-s_Ix)^{\sigma_I}}
      {T_j(x)^3T_{j^-}(-x)}.                               \tag{CF4}
\]
Here the eight charts of \((z,z,-z)\) use, respectively, charts for
\(z,-z,2z,0\): its subset sums are
\(0,z,z,-z,2z,0,0,z\).
In this order, the transition product of this diagonal cube is
\(c(z)^3c(-z)c(2z)^{-1}\).  Substitution in (CF2) cancels every
transition in (CF4).  It is nonzero, entire and lattice-periodic,
hence constant by the same argument.  Its value at zero gives
\[
 F(2x)\prod_I F(s_Ix,s_Ix,-s_Ix)^{\sigma_I}F(0)^3
                         =F(x)^3F(-x).                    \tag{CF5}
\]

Near zero take the analytic logarithm
\(f(x)=\log(F(x)/F(0))\), so \(f(0)=0\).
The face identities give \(f=0\) whenever one of the three vector
variables is zero.  The source obtains
\(f(x)=O(\|x\|^3)\) by applying the one-variable Schwarz estimate
successively to those three variables.  More explicitly, for a fixed
small nonzero \(y\), put \(\rho=R/\|y\|>1\) and consider
\(f(t_0y_0,t_1y_1,t_2y_2)\) on the polydisc
\(|t_a|<\rho\), where it is bounded by \(1\).
Divide first by \(t_0\), then by \(t_1\), then by \(t_2\), using
the face zero and the removable value at zero at each step.
The successive Schwarz estimates give at \(t_0=t_1=t_2=1\)
the bound \(\rho^{-3}=\|y\|^3/R^3\).

Set
\[
 \phi(z)=f(z,z,-z),\qquad
 (\mathcal Dk)(x)=\sum_I\sigma_I k(s_Ix).
\]
Taking logarithms in (CF5) yields, near zero,
\[
       f(2x)+\mathcal D\phi(x)-3f(x)-f(-x)=0.               \tag{CF6}
\]
There is no undetermined logarithmic multiple here: the exponential
of the left side is \(1\), while that side tends to zero.  Shrink so
its modulus is less than \(2\pi\); its possible value
\(2\pi i n\) then has \(n=0\).

Solve
\[
                    k(2z)-3k(z)-k(-z)=-\phi(z).            \tag{CF7}
\]
The source uses a convergent doubling series, not division by formal
Taylor coefficients.  For a cubic-order analytic \(h\), set
\[
 A_0(z)=h(z/2),\quad
 A_{n+1}(z)=3A_n(z/2)+A_n(-z/2),\quad k=\sum_{n\ge0}A_n.
\]
If \(|h(z)|\le C\|z\|^3\) in a ball, then
\(|A_n(z)|\le2^{-n}C\|z\|^3\) there.
The series is locally normally convergent and solves
\(k(2z)-3k(z)-k(-z)=h(z)\) by shifting its summation index.
Use \(h=-\phi\).  Then \(k(0)=0\) and \(k=O(\|z\|^3)\).

Equations (CF6)--(CF7) show that
\(G=f-\mathcal Dk\) satisfies
\(G(2x)=3G(x)+G(-x)\), and \(G=O(\|x\|^3)\).
For \(M(x)=|G(x)|+|G(-x)|\) this gives
\[
 M(x)\le4^n M(2^{-n}x)
               \le2C\|x\|^3\,2^{-n}.
\]
Let \(n\to\infty\).  Therefore \(f=\mathcal Dk\) near zero.
Define
\[
          b_i(z)=c_{ii_0}(z)e^{k(z)},\qquad \kappa=F(0)
\]
on a sufficiently small ball contained in \(W_{i_0}\).
The transition law proves \(b_i=c_{ij}b_j\); exponentiating
\(f=\mathcal Dk\) and applying (CF2) gives (CF3).
All these functions are nonzero, and \(F(0)\ne0\).
\end{proof}

\begin{srlemma}[The actual extension of the frame to all of \(V\)]
\label{lem:ftb-cube-global-frame}
The \(b_i\) extend to nonvanishing holomorphic functions on \(W_i\),
still satisfying their transition law, such that (CF3) holds on every
cube chart, for all \(x\).
\end{srlemma}

\begin{proof}
First extend a local cube identity to the whole region
\[
                    \{x:s_Ix\in B(0,r)\text{ for all }I\}.
\]
This region is convex.  On its cube-chart cover, the quotient of
\(T_j\) by the alternating product of the \(b_{j_I}\) is holomorphic
and agrees on overlaps by (CF2).  Its germ at zero is the constant
\(\kappa\); the identity theorem on this connected region makes the
quotient equal to \(\kappa\) everywhere.

One can now double the radius.  For \(w\in B(0,2r)\) put
\[
                     h(w)=(w/2,w/2,-w/2),\qquad D=\{0,1\}.
\]
The \(D\)-sum is \(w\).  Every other subset sum has norm at most
\(\|w\|/2<r\), and \(\sigma_D=-1\).
On a cube chart \(j\), solve (CF3) for the \(D\)-factor and set
\[
 b'_i(w)=c_{i,j_D}(w)\,
      \frac{\kappa\prod_{I\ne D}
                   b_{j_I}(s_Ih(w))^{\sigma_I}}{T_j(h(w))}. \tag{CF8}
\]
The denominator and every numerator factor are nonzero.
Equations (CF2) and \(c_{ij}c_{jk}=c_{ik}\) show that (CF8) is
independent of the auxiliary cube chart.  Thus these expressions
glue holomorphically; they are nonzero and have the same transition
law.  On \(B(0,r)\), (CF3) proves \(b'_i=b_i\).
The preceding convex-region argument supplies the cube identity
on the larger ball.  Iterate over the balls \(B(0,2^nr)\).
The extensions agree on overlaps, so they glue on all \(W_i\).
For any fixed cube, one sufficiently large ball contains all its
eight subset sums, proving the global form of (CF3).
\end{proof}

\begin{srproposition}[The frame and its exact Appell--Humbert data]
\label{prop:ftb-actual-appell-humbert-frame}
On the same \(e,V,\Lambda\) there are a Hermitian form \(H\), an
integral bilinear form \(E\) on \(\Lambda\), a function
\(\chi:\Lambda\to\mathbb C\) of modulus \(1\), and nonzero holomorphic
frame coordinates \(f_i\) on \(W_i\), with
\[
 \begin{split}
 f_i(z)&=c_{ij}(z)f_j(z),\\
 f_i(z)&=\chi(\lambda)
       \exp\!\left(\pi H(z,\lambda)+\tfrac{\pi}{2}H(\lambda,\lambda)\right)
       f_i(z+\lambda),\\
 E(\lambda,\mu)&=\operatorname{Im}H(\lambda,\mu).             \tag{CF9}
 \end{split}
\]
The \(f_i(0)\) equal the prescribed rigidification coordinates at
zero.  Positivity of \(H\) and perfection of \(E\) are not asserted in
this proposition; they are proved from the polarization and Weil
nondegeneracy in the next step.
\end{srproposition}

\begin{proof}
For \(\lambda\in\Lambda\) define, on any chart containing \(z\),
\[
                    \alpha_\lambda(z)=b_i(z)/b_i(z+\lambda).
\]
Periodicity of \(c_{ij}\) makes this independent of \(i\).
It is an entire nonzero function.  Compare the global cube identity
at \(x=(0,y,z)\) and at \(x+(\lambda,0,0)\).
The cube generator is periodic, and cancellation of its eight
frame factors gives
\[
       \alpha_\lambda(y+z)\alpha_\lambda(0)
                         =\alpha_\lambda(y)\alpha_\lambda(z). \tag{CF10}
\]
An entire multiplicative character \(u\) on the additive vector
space, with \(u(0)=1\), is \(\exp(du_0(z))\).
Indeed differentiation of \(u(z+w)=u(z)u(w)\) gives
\(du_w=u(w)du_0\); the derivative of
\(u(z)\exp(-du_0(z))\) is zero, and its value at zero is \(1\).
Therefore
\[
 \alpha_\lambda(z)=a(\lambda)e^{B_\lambda(z)},\qquad
                       a(\lambda)=\alpha_\lambda(0)\ne0,   \tag{CF11}
\]
where \(B_\lambda\) is complex linear.
The quotient identity
\(\alpha_{\lambda+\mu}(z)=
 \alpha_\lambda(z)\alpha_\mu(z+\lambda)\)
gives, after interchanging \(\lambda,\mu\) if necessary,
\[
 a(\lambda+\mu)=a(\lambda)a(\mu)e^{B_\lambda(\mu)},\qquad
              B_{\lambda+\mu}=B_\lambda+B_\mu.              \tag{CF12}
\]
For the second identity, equality of exponentials of complex linear
forms implies equality of the forms by differentiating at zero.

Here is the explicit normal-form calculation.  Extend the additive
map \(\lambda\mapsto B_\lambda\) real linearly from the full integral
lattice to \(V\); a lattice basis is also a real basis.
The skew relation from (CF12) is
\[
        B_\lambda(\mu)-B_\mu(\lambda)\in2\pi i\mathbb Z.     \tag{CF13}
\]
Its real part vanishes, first on the lattice and then on all of
\(V\times V\) by real linearity.
Splitting a real linear map into its complex linear and conjugate
linear parts therefore writes
\[
                   B_\lambda(z)=A(\lambda,z)+K(z,\lambda),
\]
where \(A\) is symmetric complex bilinear and \(K\) is Hermitian.
Explicitly, the splitting uses the values on the real basis vectors
\(e_j,ie_j\), with coefficients
\((F(e_j)\mp iF(ie_j))/2\); the real skew identity gives the stated
symmetries.  Put \(H=K/\pi\).
Taking imaginary parts in (CF13) shows
\(\operatorname{Im}H(\lambda,\mu)\in\mathbb Z\), and its additivity
in both variables defines \(E\).

The map
\[
 \lambda\longmapsto
 \log|a(\lambda)|-\tfrac12\operatorname{Re}A(\lambda,\lambda)
                 -\tfrac12\operatorname{Re}K(\lambda,\lambda)
\]
is additive by (CF12).  Extend it real linearly to \(V\), and write
this real linear function as \(\operatorname{Re}c(z)\) for a complex
linear \(c\).  Set
\[
 Q(z)=\tfrac12 A(z,z)+c(z),\qquad
 \chi(\lambda)=a(\lambda)
   e^{-A(\lambda,\lambda)/2-K(\lambda,\lambda)/2-c(\lambda)} .
\]
Its modulus is \(1\).  Substituting these definitions into (CF11)
gives
\[
 a(\lambda)e^{B_\lambda(z)}e^{Q(z)-Q(z+\lambda)}
       =\chi(\lambda)e^{\pi H(z,\lambda)+\pi H(\lambda,\lambda)/2}.
\]
Thus \(b_i(z)e^{Q(z)}\) has the required automorphy and transition
law.  Notice \(Q(0)=0\).

Finally choose one \(i_0\) with \(0\in W_{i_0}\), and multiply every
frame coordinate by the single nonzero scalar
\[
 \frac{\text{the prescribed rigidification coordinate in }i_0}
                  {b_{i_0}(0)e^{Q(0)}}.
\]
The algebraic rigidification transition identity makes the value in
every other origin chart correct as well.  This scalar changes
neither the exponential, the lattice, nor the automorphy in (CF9).
\end{proof}

\paragraph{Frozen-source ledger and remaining steps.}
All following files are under
\srcid{src/Development/AbelianVarietyTheory/Proofs/ToricGeometry/}
at \srcid{alpoge-source-7cbf7a0c}.
Actual evaluation of (CF1)--(CF2) is
\srcid{CubeTrivialisationOnUniformisingSpace.lean}
(\srcid{94bdcfed693818f0cdc6f7c89ad5cad13958ca40}).
The periodic face and quotient argument is
\srcid{CubeSectionFacesAndDoubling.lean}
(\srcid{e1fbe94e28161601e128fd1024f1b8bf91f012a2});
it calls the periodic entire-function lemma in
\srcid{../CovolumeBounds/EntireApplyEqOfPeriodicOfSpanReal.lean}.
The series solution to (CF7) is
\srcid{DifferentiableOnDoublingIterate.lean}
(\srcid{04ea97cf72ae242c74021e29294ee9fb81a9f6e3});
the homogeneous cubic estimate is
\srcid{EventuallyEqZeroOfIsBigOCubeOfDoubling.lean}
(\srcid{4b65b8c1681c1d8ee0924fe3afa20faf0cc6a002});
the logarithmic, Schwarz and third-difference assembly is
\srcid{SeriesApplyEqZeroOfFace.lean}
(\srcid{89bfbd686accf0eebbc48ddd5acefd5a5632b35b}).
The local coordinates in (CF3) are constructed in
\srcid{ExistsLocalFrameCompatibleCube.lean}
(\srcid{5c9dc9db6d8d8b04c16fb0643375c36c6dc22444}).
The convex-region and successive radius-doubling construction
(CF8) is \srcid{ExistsGlobalFrameOfLocal.lean}
(\srcid{3fd0ef6e101051ea7a7f5a83620d81c624974784}).
Equations (CF10)--(CF12) are
\srcid{AutomorphyFactorExpLinearOfCube.lean}
(\srcid{9e2f1f482c740af89dc5a110875fb5c07be897fc});
the explicit real-linear splitting and gauge are
\srcid{AppellHumbertNormalFormOfLinearExponent.lean}
(\srcid{1d1f093ab73ab5ad6edb72b9978e2e96148e08f0}).
The actual polarising bundle, same-exponential construction and
origin-rigidification normalization are
\srcid{ExistsAppellHumbertFrameIsRig.lean}
(\srcid{16c73490578c5cc5c0939d1bb71d450477e7be78}).
This proves exactly its Hermitian-frame output.
The later positivity, Weil/perfect-pairing and section-surjectivity
producers are still separate from this chunk.
